\documentclass[11pt,a4paper]{scrartcl}

\usepackage{../../1.\space Vienanariai\space ir\space daugianariai/manostilius_lt}

\title{Kvadratinio trinario skaidymas dauginamaisiais}
\author{Andrius Berniukevičius}

\begin{document}

\maketitle

\section{Kas yra kvadratinis trinaris?}

\begin{definition}
\vocab{Kvadratinis trinaris} -- tai daugianaris, kurį galima užrašyti pavidalu $ax^2+bx+c$, kur $x$ yra kintamasis, o $a, b, c$ -- skaičiai (koeficientai) ir $a \neq 0$.
\end{definition}

Kvadratinis trinaris labai panašus į kvadratinę lygtį, tačiau tai nėra lygtis (nėra lygybės ženklo). Kad išskaidytume kvadratinį trinarį dauginamaisiais, pirmiausia turime jį prilyginti nuliui ir išspręsti gautą kvadratinę lygtį $ax^2+bx+c=0$. Šios lygties sprendiniai $x_1$ ir $x_2$ vadinami \textbf{kvadratinio trinario šaknimis}.

\section{Skaidymo dauginamaisiais formulė ir jos įrodymas}

\begin{property}[Kvadratinio trinario skaidymo formulė]
Jei kvadratinio trinario $ax^2+bx+c$ šaknys yra $x_1$ ir $x_2$, tai šį trinarį galima išskaidyti dauginamaisiais pagal formulę:
\[ ax^2+bx+c = a(x-x_1)(x-x_2). \]
\end{property}

Įrodykime šią formulę pasinaudodami Vijeto teorema.

\begin{proof}
Pirmiausia iš reiškinio $ax^2+bx+c$ iškelkime koeficientą $a$ prieš skliaustus:
\[ ax^2+bx+c = a\left(x^2 + \frac{b}{a}x + \frac{c}{a}\right). \]
Prisiminkime Vijeto teoremą, kuri sieja lygties šaknis su jos koeficientais:
\[
\begin{aligned}
x_1 + x_2 &= -\frac{b}{a}, \quad \text{todėl} \quad \frac{b}{a} = -(x_1 + x_2), \\
x_1 \cdot x_2 &= \frac{c}{a}.
\end{aligned}
\]
Įstatykime šias išraiškas į mūsų skliaustuose esantį reiškinį vietoje $\frac{b}{a}$ ir $\frac{c}{a}$:
\[ a\left(x^2 - (x_1 + x_2)x + x_1 x_2\right). \]
Atskliauskime vidinius skliaustus:
\[ a\left(x^2 - x_1 x - x_2 x + x_1 x_2\right). \]
Dabar pritaikykime grupavimo būdą -- sugrupuokime pirmus du narius ir paskutinius du narius. Iš pirmų dviejų iškeliame $x$, o iš kitų dviejų iškeliame $-x_2$:
\[
\begin{aligned}
a\left( x(x - x_1) - x_2(x - x_1) \right).
\end{aligned}
\]
Iškeliame bendrą skliaustą $(x - x_1)$:
\[ a(x - x_1)(x - x_2). \]
Įrodėme, kad $ax^2+bx+c = a(x-x_1)(x-x_2)$.
\end{proof}

\begin{remark}
Jei kvadratinio trinario diskriminantas $D = 0$, jis turi vieną šaknį $x_1$ (arba dvi sutampančias šaknis $x_1 = x_2$). Tuomet formulė atrodo taip: $a(x-x_1)^2$.
Jei diskriminantas $D < 0$, trinaris realiųjų šaknų neturi ir jo išskaidyti realiaisiais daugikliais \textbf{negalima}.
\end{remark}

\section{Pavyzdžiai}

\begin{example}
Išskaidykite dauginamaisiais kvadratinį trinarį:
\[ x^2 - 7x + 10 \]
\end{example}
\begin{solution}
Prilyginame trinarį nuliui, kad rastume jo šaknis:
\[ x^2 - 7x + 10 = 0 \]
Pagal Vijeto teoremą (arba skaičiuojant diskriminantą) randame šaknis:
\[ x_1 = 2, \quad x_2 = 5. \]
Kadangi koeficientas $a=1$, taikome formulę $a(x-x_1)(x-x_2)$:
\[ 1 \cdot (x - 2)(x - 5) = (x - 2)(x - 5). \]
\textbf{Atsakymas:} $(x - 2)(x - 5)$.
\end{solution}

\begin{example}
Išskaidykite dauginamaisiais kvadratinį trinarį:
\[ -2x^2 + 10x - 12 \]
\end{example}
\begin{solution}
Prilyginame nuliui ir randame šaknis:
\[ -2x^2 + 10x - 12 = 0 \]
Padaliję iš $-2$, gauname $x^2 - 5x + 6 = 0$. Šios lygties šaknys yra $x_1 = 2$ ir $x_2 = 3$.
Pradinio trinario koeficientas $a = -2$. Įstatome viską į formulę:
\[ -2(x - 2)(x - 3). \]
\textbf{Atsakymas:} $-2(x - 2)(x - 3)$.
\end{solution}

\begin{example}
Išskaidykite dauginamaisiais kvadratinį trinarį:
\[ 3x^2 + 5x - 2 \]
\end{example}
\begin{solution}
Sprendžiame lygtį $3x^2 + 5x - 2 = 0$. Skaičiuojame diskriminantą:
\[ D = 5^2 - 4 \cdot 3 \cdot (-2) = 25 + 24 = 49. \]
\[ x_1 = \frac{-5 - 7}{6} = -2, \quad x_2 = \frac{-5 + 7}{6} = \frac{2}{6} = \frac{1}{3}. \]
Taikome skaidymo formulę su $a=3$:
\[ 3(x - (-2))\left(x - \frac{1}{3}\right) = 3(x + 2)\left(x - \frac{1}{3}\right). \]
Kad atsakymas atrodytų gražiau ir neturėtų trupmenų, daugiklį $3$ galime įkelti į antruosius skliaustus:
\[ (x + 2)\left(3 \cdot x - 3 \cdot \frac{1}{3}\right) = (x + 2)(3x - 1). \]
\textbf{Atsakymas:} $(x + 2)(3x - 1)$.
\end{solution}

\section{Uždaviniai}

\begin{problem}
Išskaidykite dauginamaisiais šiuos kvadratinius trinarius (kai $a=1$):
\begin{tasks}[label={(\arabic*)}, label-offset=0.9em](2)
    \task $x^2 - 6x + 8$
    \task $x^2 + 4x - 5$
    \task $x^2 - x - 12$
    \task $x^2 + 8x + 15$
    \task $x^2 - 9x + 14$
    \task $x^2 + 3x - 18$
    \task $2x^2 - 14x + 24$
    \task $-3x^2 - 6x + 9$
    \task $2x^2 + 5x - 3$
    \task $5x^2 - 3x - 2$
    \task $3x^2 - 7x - 6$
    \task $-2x^2 + x + 3$
\end{tasks}
\end{problem}

\begin{problem}
Išskaidykite kvadratinį trinarį ir suprastinkite algebrines trupmenas:
\begin{tasks}[label={(\arabic*)}, label-offset=0.9em](2)
    \task $\frac{x^2 - 3x - 10}{x - 5}$
    \task $\frac{x^2 - 9}{x^2 - x - 6}$
    \task $\frac{2x^2 + x - 1}{x + 1}$
    \task $\frac{3x^2 - 12x}{x^2 - 2x - 8}$
    \task $\frac{x^2 - 4x - 12}{x - 6}$
    \task $\frac{x^2 - 4}{x^2 + x - 6}$
\end{tasks}
\end{problem}

\newpage
\answerssection

\begin{problem} \phantom{text}
\begin{tasks}[label={(\arabic*)}, label-offset=0.9em](2)
    \task $(x-2)(x-4)$;
    \task $(x-1)(x+5)$;
    \task $(x+3)(x-4)$;
    \task $(x+3)(x+5)$;
    \task $(x-2)(x-7)$;
    \task $(x-3)(x+6)$.
    \task $2(x-3)(x-4)$;
    \task $-3(x-1)(x+3)$;
    \task $(x+3)(2x-1)$;
    \task $(x-1)(5x+2)$;
    \task $(3x+2)(x-3)$;
    \task $(3-2x)(x+1)$.
\end{tasks}
\end{problem}

\begin{problem} \phantom{text}
\begin{tasks}[label={(\arabic*)}, label-offset=0.9em](2)
    \task $x+2$;
    \task $\frac{x-3}{x-2}$;
    \task $2x-1$;
    \task $\frac{3x}{x+2}$;
    \task $x+2$;
    \task $\frac{x+2}{x+3}$.
\end{tasks}
\end{problem}

\end{document}